\documentclass[10pt,a4paper]{article}
\usepackage[margin=0.85in]{geometry}
\usepackage{booktabs,amsmath,amssymb,graphicx,microtype}
\usepackage[table]{xcolor}
\usepackage[colorlinks=true,linkcolor=black,citecolor=black,urlcolor=blue]{hyperref}
\usepackage{titlesec}
\titlespacing*{\section}{0pt}{6pt}{3pt}
\titlespacing*{\subsection}{0pt}{5pt}{2pt}
\titleformat{\section}{\large\bfseries}{\thesection.}{0.5em}{}
\titleformat{\subsection}{\normalsize\bfseries}{\thesubsection}{0.5em}{}
\setlength{\parskip}{2pt}
\newcommand{\num}[1]{\textbf{#1}}

\begin{document}
\pagestyle{empty}

\begin{center}
{\Large\bfseries A Quantum Allosteric Scanner for Undruggable Targets}\\[2pt]
{\large Continuous-time quantum walks with learned operator selection}\\[4pt]
{\normalsize Team AuraQu \quad$\vert$\quad Cleveland Clinic problem statement:\\
\emph{Unlocking undruggable targets: quantum simulation of allosteric signal propagation}}
\end{center}
\vspace{-4pt}\hrule\vspace{6pt}

\section{Problem Framing}

A cryptic allosteric pocket is useful precisely when it is \emph{distal}: far enough from the
active site that no geometric pocket-finder looking near the catalytic machinery will locate it.
This is not a rhetorical point, and we measured it. Ranking residues purely by graph distance from
the active site --- the strongest ``free'' classical heuristic --- achieves \num{AUC 0.227} on
distal targets across 53 protein families. That is not merely weak; it is \emph{below chance},
because a distal pocket is by definition where proximity is not. Geometry does not underperform on
this problem, it points the wrong way.

That is the gap a quantum method can fill, and it is why the challenge's own success definition
names \emph{distal} regulatory residues specifically. Signal propagation through a protein contact
network is a transport problem on a graph. A continuous-time quantum walk (CTQW) explores such a
graph by amplitude interference rather than diffusion: paths cancel and reinforce, so the walker
distinguishes routes that a classical random walk --- which only accumulates probability ---
cannot. Our hypothesis is that this interference structure is what carries allosteric information,
and our evidence is that the walk succeeds exactly where the distance heuristic inverts.

\section{Technical Approach}

\textbf{Paradigm.} Hybrid quantum--AI. The physics layer is a CTQW, currently executed by exact
spectral simulation and targeted for Trotterised circuit execution in Phase 2 (\S3). The learning
layer is a classical regressor that selects which quantum operator to run. We state plainly that
today's results come from simulation, not hardware; \S3 gives the qubit and depth budget and the
path to execution.

\textbf{The walk.} From the apo structure alone we build a C$\alpha$ contact graph at 8\,\AA{}
(elastic-network assumption, as the challenge specifies) and a Hamiltonian $H$ from its adjacency
or Laplacian, optionally with a diagonal site potential $V$ encoding B-factor disorder, solvent
exposure, rigidity and low-mode participation. Seeding at the active site, the walker evolves as
$U(t)=e^{-iHt}$ and we score every candidate residue by its dynamic connectivity to that site.

A single evolution time is unsafe: destructive interference can drive a genuine coupling to zero,
and on disordered graphs the walk can Anderson-localise. We therefore use time-averaged
quantities, including the closed form
$p_{\mathrm{avg}}(s\!\to\!a)=\sum_k |v_k(s)|^2|v_k(a)|^2$,
which is parameter-free and provably free of cancellation. Alongside it we score Green's functions
$|G(E,\eta)|^2$, walker dispersion diagnostics, and a spectral variance
$\Delta E=\sqrt{\langle E^2\rangle-\langle E\rangle^2}$ that proved to be our single strongest
score.

\textbf{The operator is not universal --- and that is the central technical finding.} We evaluated
13 Hamiltonians $\times$ 17 scores $\times$ 4 seed-filter settings $=$ 884 configurations per
protein on 630 proteins. A variance decomposition is decisive: main effects (``this Hamiltonian is
generally good'') explain \num{1\%} of the variance in performance; the remaining \num{99\%} is
interaction and per-protein variation. Across 630 proteins there are \num{308 distinct winning
configurations}. There is no single right operator to publish.

\textbf{The AI layer follows directly from that.} Rather than fix an operator or search all 884,
we train a model to \emph{recommend} configurations from protein topology alone. Input: nine graph
features (size, edges, degree statistics, clustering, diameter, algebraic connectivity, spectral
radius). Output: a ranked, weighted shortlist of $k$ configurations to run, $k$ chosen by the user.
A feasibility gate first removes settings the protein cannot support --- on a compact protein
spanning five hops, an aggressive seed filter deletes the answer before the walk starts.

\section{Feasibility and Resource Requirements}

\textbf{Already demonstrated.} The pipeline described above has been run end to end on a
1022-protein worklist; 630 proteins across 399 families produce a scoreable result. Total cloud
compute for every experiment reported here is under \num{EUR 3}. The classical simulation cost
is dominated by one eigendecomposition per operator per protein, which is seconds at these sizes.
Feasibility of the \emph{classical} half is therefore not a projection --- it is measured.

\textbf{The quantum half, honestly.} Mapping one residue to one qubit is not viable for a
500-residue target on near-term hardware. Our Phase 2 plan is basis encoding of the seed with
Trotterised evolution, $e^{-iHt}\approx(\prod_j e^{-iH_jt/r})^r$, where $H$ decomposes over graph
edges. At $r=2$ we already measure $\sim10^{-3}$ agreement with exact evolution in simulation, so
depth stays shallow by construction. Coarse-graining is the enabling step and we report a genuine
failure: a 16-node community-level walk collapsed to AUC 0.529, because scoring at community level
and back-propagating to residues forces ties. Fixing this --- compressing while proving the
topological signal survives --- is the principal Phase 2 risk and the principal Phase 2 task.

\textbf{Resources needed.} AWS Braket and Classiq access as provided; one modest CPU node for the
classical layer. No proprietary data: every structure is public RCSB, and no classical MD
trajectory enters the pipeline, per the challenge constraint.

\section{Expected Impact}

A successful PoC delivers a scanner that ingests an apo structure of a shelved target and returns a
ranked probability map of candidate allosteric sites, upstream of high-throughput screening. The
quantitative target we propose for Phase 2 follows from what we already achieve: on distal targets
the walk must beat the geometric floor with family-level significance, which it currently does at
$p<10^{-8}$, and it must place true pocket residues in the top five often enough to direct library
design. We are explicit that the second half is where the current method is weakest (\S5), and that
is the honest measure of what Phase 2 must improve.

\section{Validation Plan}

Our validation discipline is the part of this proposal we would most want a reviewer to examine,
because the failure mode in this field is a number that survives no control.

\textbf{Counting.} Structures are grouped into protein families and every result is counted per
family, never per structure. This is not cosmetic: one family in our set contributes 28 structures
of a single protein, and per-structure counting inflated an early result by roughly twentyfold.

\textbf{Controls actually run.} Label-permutation nulls preserving score correlation; family-level
permutation nulls; a favourite-pocket null; a random-residue-through-the-same-pipeline null; and
the proximity floor described in \S1. Classical baselines --- proximity, degree, pocket size,
fpocket druggability, PASSer rank --- are scored on the identical residues and labels.

\textbf{Blind evaluation.} Every model number is leave-one-family-out cross-validated: each protein
is scored by a model that never saw its family.

\textbf{Headline results, and their limits.}

\begin{center}\small
\begin{tabular}{lrrr}
\toprule
& proximity floor & CTQW (blind) & families where CTQW wins \\ \midrule
\textbf{Distal targets} (91 proteins, 53 families) & 0.227 & \num{0.617} & \num{47/53}, $p<10^{-8}$ \\
Near targets (539 proteins, 350 families) & 0.633 & 0.597 & 137/350 (floor wins) \\
All (630 proteins, 399 families) & 0.574 & 0.600 & 182/399, $p=0.09$ \\
\bottomrule
\end{tabular}
\end{center}

We state the negative half plainly: \emph{overall} the walk does not beat distance. The result is
specifically and only about distal targets --- which is what the challenge asks about, and which
is the case classical geometry cannot address.

On held-out targets never seen in training, the model's six recommended configurations were run
through the full pipeline. On \textbf{BCR-ABL1} (1OPL$\to$5MO4, asciminib) the best of the six
reached \num{AUC 0.900, P@5 0.80} --- four of the top five residues are genuine myristoyl-pocket
contacts. On \textbf{HIV-1 reverse transcriptase} the best reached AUC 0.679, matching an
independent pre-registered run, but with no top-five hit.

\textbf{What we found wrong, and report anyway.} The recommender's \emph{ranking} works; its
\emph{confidence weighting} does not --- on BCR-ABL1 the winning configuration sat at rank 4 with
7\% weight while the top-weighted pick scored AUC 0.140. Merging the shortlist is worse than
picking one from it. And two defects in the challenge's own Table 1 warrant flagging to the
organisers: \textbf{6C1H contains no mavacamten} (only ADP and MG, verified live against RCSB), and
\textbf{4OBE is wild-type KRAS}, not G12C --- residue 12 is glycine, so the named apo structure is
the wrong genotype for the stated objective.

\section{Hybrid / Cross-Domain Integration}

The hybrid design is not decorative. The quantum layer produces evidence --- a residue ranking per
operator --- and the AI layer decides which evidence to gather, reducing 884 candidate quantum
computations to $k$ (typically 6), a $147\times$ reduction. Measured leave-one-family-out, $k{=}6$
retains a configuration reaching AUC $\ge0.8$ for 45\% of proteins and recovers 43 of the 71
families that exhaustive search finds; $k{=}20$ recovers 54. This matters more, not less, on
quantum hardware, where each configuration is a separate circuit and exhaustive sweeps are
impossible. HPC orchestration runs the classical sweeps that generate the model's training data.

\section{Team Capability}

\textcolor{red}{[PLACEHOLDER --- team names, roles, affiliations, and prior quantum/domain
experience to be supplied. Note the guidelines require a separate Team Profile section in the
portal in addition to this.]}

\end{document}
